% !TEX program = xelatex
% Рабочий черновик обзорной статьи. Неизвестные сведения не заполняются вымышленными данными.

\documentclass{misisarticle}

\grnti{28.23.27}
\asjc{1702, 2207}
\oecd{2.2}

\title{Обзор методов тактического планирования и командного управления в RoboCup Small Size League}

\misisauthor{С.~Л.~Трафняк}{кафедра инженерной кибернетики НИТУ МИСИС, группа БПМ-23-РК-1}{Москва, Россия}{}

\abstract{%
  Рассматриваются методы принятия тактических решений в международных командах RoboCup Small Size League. Цель обзора~--- систематизировать подходы к распределению ролей, выбору атакующих действий, организации передач, защите и учёту соперника, а также определить решения, применимые при разработке стратегии собственной команды. Анализ построен по решаемым задачам, а не как последовательное описание команд; основными источниками служат публикации команд и доступные версии программного кода. Показано, что современные стратегии сочетают иерархическое разбиение поведения, оценочные функции, прогноз достижимости и механизмы стабилизации решений. Обучаемые многоагентные подходы расширяют возможности адаптации, однако пока требуют решения вопросов соблюдения правил и переноса из симуляции на реальные роботы. Результатом обзора является сравнительная схема методов и набор требований к будущему алгоритму командного управления на базе инфраструктуры PAcmaCS.
}

\keywords{%
  RoboCup Small Size League, многоагентные робототехнические системы, тактическое планирование, распределение ролей, выбор передачи, моделирование соперника
}

\begin{document}

\maketitle

\section{Введение}

RoboCup Small Size League (SSL) представляет собой динамическую состязательную среду, в которой автономная команда колёсных роботов должна в реальном времени интерпретировать состояние поля, распределять функции между роботами и согласованно выполнять игровые действия. Глобальная система технического зрения предоставляет координаты объектов, а программа команды принимает решения о ролях, перемещениях, передачах, ударах и защитных действиях. Поэтому качество игры определяется не только точностью отдельных навыков, но и тем, как они объединены в общую тактическую систему.

Классической основой такого объединения является архитектура STP (Skills, Tactics and Plays), разделяющая низкоуровневые навыки, индивидуальные тактики и командные розыгрыши~\cite{stp2005}. Однако современные команды развивают эту схему по-разному: применяют приоритетное распределение ролей, агентные переговоры, генерацию и многоступенчатую оценку передач, прогноз времени овладения мячом и обучаемые политики. Сопоставление этих решений затруднено тем, что командные публикации относятся к разным годам, описывают разные подзадачи и не всегда содержат полную турнирную стратегию.

Цель работы~--- систематизировать опубликованные методы тактического планирования команд RoboCup SSL и оценить их применимость при разработке собственного алгоритма управления командой роботов. Исследовательский вопрос формулируется следующим образом: как команды SSL принимают тактические решения, чем различаются их подходы и какие решения могут быть перенесены в стратегию, создаваемую на основе инфраструктуры SPbUnited PAcmaCS~\cite{pacmacs}.

В обзор включены задачи выбора действий команды, назначения ролей, поиска позиций поддержки, организации передач и ударов, защиты и учёта соперника. Вопросы конструкции роботов, радиосвязи и технического зрения рассматриваются только как ограничения, влияющие на исполнимость тактического решения.

\section{Методика обзора}

Работа является целевым аналитическим обзором, а не систематическим обзором литературы. Основной корпус составили статьи STP и Team Description Papers команд SPbUnited, TIGERs Mannheim, ZJUNlict, ER-Force, RobôCIn и RoboJackets за 2019--2025~гг.~\cite{spb2025,tigers2025,zjunlict2025,erforce2019,robocin2025,robojackets2024}. Для оценки перспектив обучаемого управления дополнительно рассмотрена работа TRAPS 2026~г. по многоагентному обучению с подкреплением~\cite{traps2026}. Там, где это возможно, описание сопоставлялось с открытым программным кодом, в частности с реализацией оценки перехвата передачи TIGERs~\cite{tigerscode}.

Единицей сравнения является не команда в целом, а отдельный алгоритмический компонент. Для каждого компонента фиксируются решаемая задача, входные данные, способ формирования решения, учёт соперника, способ проверки и ограничения. Такое разбиение не позволяет ошибочно сравнивать, например, генератор точек передачи с алгоритмом управления приёмом: эти компоненты могут быть последовательными частями одной системы.

При интерпретации соблюдаются следующие ограничения. Описание в публикации не считается доказательством использования той же версии алгоритма в текущем турнире. Наличие открытого framework не означает публикацию полной стратегии. Турнирный результат сам по себе не доказывает превосходство конкретного метода. В настоящей работе не выполнялась самостоятельная количественная обработка матчевых логов, поэтому выводы относятся к опубликованным алгоритмам, а не к статистически измеренной эффективности команд.

\section{Архитектуры принятия решений}

\subsection{Иерархическое управление}

STP задаёт удобную декомпозицию командного поведения. Навык описывает исполнимое действие одного робота, тактика выбирает и параметризует последовательность навыков, а розыгрыш координирует тактики нескольких роботов~\cite{stp2005}. Преимущество подхода состоит в разделении временных масштабов: робот может быстро реагировать на изменение обстановки, не отказываясь от более продолжительной командной цели. Недостатком является рост числа правил и переходов при усложнении стратегии.

ZJUNlict сохраняет иерархическую организацию, но явно разделяет глобальное планирование, индивидуальное принятие решений и многоагентное исполнение~\cite{zjunlict2025}. Глобальный уровень определяет состояние атаки или защиты и выбирает тактическую схему; следующий уровень формирует навык отдельного робота; нижний уровень выполняет распределение задач и динамическое сопоставление ролей. Механизм удержания состояния уменьшает частые переключения между близкими решениями. Это существенно для SSL, поскольку небольшое изменение координат не должно немедленно отменять уже начатую передачу или удар.

\subsection{Распределение ролей и агентные подходы}

В SPbUnited каждой роли сопоставлены приоритет, метрика выбора исполнителя и собственное поведение~\cite{spb2025}. Сначала наиболее приоритетная роль выбирает подходящего робота, после чего остальные роли распределяются между оставшимися роботами. Метрика может учитывать расстояние до целевой позиции и дополнительные условия, зависящие от роли. Такой алгоритм прозрачен и вычислительно прост, однако последовательный выбор не гарантирует глобально оптимального назначения всего набора ролей.

RoboJackets описывает переход от централизованного ``тренера'' к агентной архитектуре~\cite{robojackets2024}. Каждый агент самостоятельно выбирает своё поведение, а необходимые согласования выполняются через межагентный обмен. В атаке владеющий мячом робот рассылает запрос на передачу, потенциальные получатели оценивают собственную доступность, после чего выбирается действие. Подход уменьшает зависимость от единственного координирующего модуля, но повышает требования к разрешению конфликтов и согласованности локальных решений.

\section{Атакующие действия}

\subsection{Генерация и выбор передач}

SPbUnited формулирует поиск точки передачи как оптимизацию оценочной функции. Полный перебор поля дорог, а вычисление градиента не всегда удобно, поэтому локальный поиск выполняется методом Нелдера--Мида из ранее найденных и случайно созданных начальных точек~\cite{spb2025}. Метод позволяет получить достаточно хорошее решение за приемлемое время, но не гарантирует нахождения глобального оптимума. Качество результата полностью зависит от состава и настройки оценочной функции.

Наиболее подробно конвейер работы с передачами описан TIGERs Mannheim~\cite{tigers2025}. Сначала для достижимых позиций формируются варианты прямой и навесной передачи с параметрами времени и скорости. Затем исключаются варианты, нарушающие правила или геометрические ограничения. Оставшиеся передачи независимо оцениваются по риску перехвата, возможности последующего удара, продвижению и создаваемому давлению. Выбор выполняется по последовательности приоритетов: передача под удар в касание, передача с последующим ударом, безопасное продвижение и запасной вариант. Раздельные оценки облегчают объяснение решения и не скрывают компромисс в одной суммарной величине.

Риск перехвата у TIGERs оценивается по достижимым областям соперников вдоль траектории мяча~\cite{tigers2025,tigerscode}. Для каждого момента движения мяча рассматриваются области, которых противник может достичь при разных динамических предположениях. Плавная оценка предпочтительнее жёсткого порога, поскольку снижает нестабильное переключение между почти равными передачами. Дополнительная адаптация изменяет модель движения соперников с учётом фактической успешности передач; заявленная командой цель в 70\% является параметром их настройки, а не универсальным нормативом.

\subsection{Поддержка и приём мяча}

ER-Force решает задачу продвижения мяча через промежуточные позиции полузащитников~\cite{erforce2019}. Кандидатные позиции оцениваются по близости противников, продвижению к воротам, длине передачи, возможности продолжить атаку без остановки мяча и времени прибытия получателя. Роботам назначаются непересекающиеся области поиска, что уменьшает конкуренцию за одну позицию и формирует цепочку коротких передач.

RobôCIn сосредоточен на согласовании траектории мяча и движения получателя~\cite{robocin2025}. Алгоритм ищет раннюю точку, в которой робот успевает встретить движущийся мяч с заданным временным запасом. Для ускорения применяется двоичный поиск по траектории, а целевая точка дополнительно стабилизируется. Это показывает, что хороший выбор направления передачи недостаточен: стратегия должна проверять динамическую достижимость приёма и обеспечивать устойчивость целевой позиции.

\section{Защита и учёт соперника}

Учёт соперника в рассмотренных системах реализуется на разных уровнях. В простейшем случае его текущие координаты входят в геометрические ограничения передачи или удара. Более содержательный вариант использует прогноз достижимости: модель отвечает не только на вопрос, где находится соперник, но и успеет ли он к определённой точке к моменту прохождения мяча.

ZJUNlict прогнозирует время перехвата или преследования свободного мяча для роботов обеих команд, после чего сравнивает минимальные времена~\cite{zjunlict2025}. Результат используется как прогноз будущего владения и влияет на выбор агрессивного либо консервативного режима. Такой прогноз описывает физическую достижимость и не должен интерпретироваться как распознавание намерений соперника.

ER-Force использует оценку перспективности получателя передачи для защитной опеки~\cite{erforce2019}. Из рассмотрения исключаются малозначимые цели, а оставшиеся соперники оцениваются по удалённости от мяча и собственных ворот. После назначения роботов для блокирования наиболее опасных ударных позиций свободные защитники получают задачи опеки возможных адресатов передачи. Тем самым защита сочетает предотвращение непосредственного удара и попытку перехватить развитие атаки.

В TIGERs модель риска обновляется на основе успешности передач в целом и между зонами поля~\cite{tigers2025}. Этот механизм можно рассматривать как ограниченную онлайн-адаптацию параметров, но не как полное распознавание тактики соперника. Более сложное моделирование возможно на основе последовательностей действий: работа RoboDragons демонстрирует онлайн-кластеризацию розыгрышей по извлечённым действиям~\cite{robodragons2019}. Однако исследовательский результат не подтверждает использование соответствующего алгоритма в актуальной турнирной стратегии.

\section{Обучаемые подходы}

Обучение с подкреплением потенциально уменьшает объём ручного проектирования правил, но усложняет интерпретацию, гарантии соблюдения ограничений и перенос на реальную аппаратуру. TRAPS использует многоагентный вариант PPO (MAPPO), централизованное обучение с децентрализованным исполнением, self-play и поэтапную программу обучения в Isaac Lab~\cite{traps2026}. В экспериментах 3$\times$3 наблюдалось формирование совместных действий, в том числе подстраховки партнёра.

Полученный результат пока нельзя считать готовой заменой турнирной стратегии. Авторы отмечают нарушения правил, в частности заезд в штрафную область, а полноценный перенос политики на физические матчи 3$\times$3 относят к дальнейшей работе~\cite{traps2026}. Поэтому практически оправдан гибридный вариант: обучаемые компоненты применяются для ограниченных подзадач или оценки ситуации, а жёсткие правила безопасности и регламентные ограничения остаются в детерминированном слое.

\section{Сравнительный анализ}

В табл.~\ref{tab:comparison} методы сопоставлены по основной задаче, механизму решения и ключевому ограничению. Таблица отражает опубликованные компоненты, а не полную силу команд и не их турнирный рейтинг.

\begin{table*}[t]
  \caption{Сравнение рассмотренных компонентов тактического управления}
  \label{tab:comparison}
  \centering
  \small
  \begin{tabular}{@{}p{0.11\textwidth}p{0.16\textwidth}p{0.29\textwidth}p{0.27\textwidth}@{}}
    \toprule
    \textbf{Источник} & \textbf{Задача} & \textbf{Подход} & \textbf{Ограничение} \\
    \midrule
    SPbUnited & Роли и точки передач & Приоритетное назначение ролей; локальная оптимизация Нелдера--Мида & Последовательное назначение и локальный поиск не гарантируют глобальный оптимум \\
    TIGERs & Выбор передачи & Генерация, фильтрация, раздельное оценивание и поэтапный выбор & Требуются точные модели движения и настройка множества порогов \\
    ZJUNlict & Владение и иерархия решений & Сравнение времени достижения мяча; удержание состояния & Прогноз достижимости не описывает намерения соперника \\
    ER-Force & Поддержка и защита & Области поиска полузащитников; оценка адресатов и опека & Тактическое описание относится к 2019~г. \\
    RobôCIn & Приём движущегося мяча & Двоичный поиск ранней достижимой точки с временным запасом & Решается локальная подзадача, а не выбор командного действия \\
    TRAPS & Обучаемая кооперация & MAPPO, self-play, curriculum, симуляция 3$\times$3 & Соблюдение правил и sim-to-real остаются открытыми задачами \\
    \bottomrule
  \end{tabular}
\end{table*}

Для собственной системы на базе PAcmaCS целесообразно сохранить явную иерархию и модульность, поскольку они упрощают отладку в симуляторе и последующий перенос на роботов. В качестве начальной архитектуры предлагается: приоритетное распределение обязательных ролей; независимая генерация кандидатов действий; отдельные оценки достижимости, риска перехвата и тактической ценности; механизм гистерезиса или удержания состояния; журналирование причин выбора. Учёт соперника следует начинать с прогноза времени достижения и риска перехвата, не смешивая его с более сложным распознаванием намерений.

Сравнение выявляет также требования к будущей проверке алгоритма. Необходимо фиксировать версии исходного кода и параметры моделей, разделять результаты симуляции и реальных роботов, а для каждой подзадачи заранее определить метрику. Для передач такими метриками могут быть доля успешных передач, частота перехватов и продвижение мяча; для распределения ролей~--- время перестроения и число конфликтующих назначений; для устойчивости решений~--- частота необоснованных переключений. Конкретная экспериментальная программа выходит за рамки настоящего обзора и должна быть сформирована после выбора реализуемого алгоритма.

\section{Заключение}

Опубликованные системы RoboCup SSL показывают, что тактическое управление обычно строится как композиция специализированных компонентов, а не как единый универсальный алгоритм. Иерархические архитектуры обеспечивают понятную декомпозицию; приоритетные и агентные схемы решают задачу распределения ответственности; конвейеры передач соединяют генерацию, проверку достижимости, оценку риска и выбор действия; прогноз времени прибытия позволяет учитывать соперника без предположений о его намерениях.

Для разработки собственной стратегии наиболее применима модульная гибридная архитектура с объяснимыми оценками и устойчивостью к частым переключениям. Обучаемые компоненты перспективны, однако на текущем этапе не устраняют необходимость детерминированного контроля правил и проверки на реальной аппаратуре. Следующий этап работы~--- детализировать карточки алгоритмов выбранных команд, проверить соответствие описаний доступному коду и сформировать спецификацию первого реализуемого прототипа для PAcmaCS.

\begin{thebibliography}{99}

\bibitem{stp2005} Browning B., Bruce J., Bowling M., Veloso M. STP: Skills, Tactics and Plays for Multi-Robot Control in Adversarial Environments. 2005. URL: \url{https://www.cs.cmu.edu/~mmv/papers/05STP-IEEE.pdf} (дата обращения: 29.10.2026).

\bibitem{spb2025} Fradkov A., Filippov S., Meshcheryakov A. et al. SPbUnited TDP for RoboCup 2025. 2025. URL: \url{https://ssl.robocup.org/wp-content/uploads/2025/04/2025_TDPSPbUnited.pdf} (дата обращения: 29.10.2026).

\bibitem{pacmacs} SPbUnited. PAcmaCS: программная инфраструктура команды. URL: \url{https://github.com/SPBUnited/PAcmaCS} (дата обращения: 29.10.2026).

\bibitem{tigers2025} Geiger M., Ommer N., Ryll A., Ratzel M. TIGERs Mannheim Extended Team Description for RoboCup 2025. 2025. URL: \url{https://ssl.robocup.org/wp-content/uploads/2025/04/2025_ETDP_TIGERs-Mannheim.pdf} (дата обращения: 29.10.2026).

\bibitem{tigerscode} TIGERs Mannheim. Реализация динамической оценки перехвата передачи. \href{https://github.com/TIGERs-Mannheim/Sumatra/blob/496376d77d9e208736257f4410ad1dd5ca9cf72d/modules/moduli-ai/src/main/java/edu/tigers/sumatra/ai/metis/pass/rating/DynamicPassInterceptionMovingRobotRater.java}{Код в GitHub, commit 496376d} (дата обращения: 29.10.2026).

\bibitem{zjunlict2025} Wu Z., Wang L., Yang Z. et al. ZJUNlict Extended Team Description Paper. Small Size League of RoboCup 2025. 2025. URL: \url{https://ssl.robocup.org/wp-content/uploads/2025/04/2025_ETDP_ZJUNlict.pdf} (дата обращения: 29.10.2026).

\bibitem{erforce2019} Engelhardt T., Heineken T., Kuehn T. et al. ER-Force 2019 Extended Team Description Paper. 2019. URL: \url{https://ssl.robocup.org/wp-content/uploads/2019/03/2019_ETDP_ER-Force.pdf} (дата обращения: 29.10.2026).

\bibitem{robocin2025} Franca A., Barros B., Barbosa D. C. et al. RobôCIn Small Size League Extended Team Description Paper for RoboCup 2025. 2025. URL: \url{https://ssl.robocup.org/wp-content/uploads/2025/04/2025_ETDP_RoboCIn.pdf} (дата обращения: 29.10.2026).

\bibitem{robojackets2024} Nayak P., Das M., Parikh S. et al. RoboJackets 2024 Team Description Paper. 2024. URL: \url{https://ssl.robocup.org/wp-content/uploads/2024/04/2024_TDP_RoboJackets.pdf} (дата обращения: 29.10.2026).

\bibitem{robodragons2019} Adachi K., Ito Y., Naruse T. Online Strategy Clustering Based on Action Sequences in RoboCupSoccer Small Size League. Robotics. 2019. Vol.~8, no.~3. URL: \url{https://ssl.robocup.org/wp-content/uploads/2022/11/robotics-08-00058.pdf} (дата обращения: 29.10.2026).

\bibitem{traps2026} Naito Y., Mitsuoka R., Inagaki K. TRAPS Extended Team Description for RoboCup 2026. 2026. URL: \url{https://ssl.robocup.org/wp-content/uploads/2026/03/2026_TDP_TRAPS.pdf} (дата обращения: 29.10.2026).

\end{thebibliography}

\end{document}
